% !TeX document-id = {ecbf6175-be07-4601-b8eb-d3263d2bff07}
% !TeX TXS-program:compile = txs:///pdflatex/[--shell-escape]

%include guard:
\ifdefined\GlobalAlreadyIncluded
  \expandafter\endinput
\fi

\gdef\GlobalAlreadyIncluded{}
%

\documentclass[accentcolor=3c,landscape,ngerman,presentation,t,usenames,dvipsnames,svgnames,table, aspectratio=169]{tudabeamer}

% Template-Modifikationen
\addtobeamertemplate{frametitle}{}{\vspace{-1em}} % mehr Platz vor dem Inhalt

% andere global gemeinsame definitionen
%Includes
\usepackage[ngerman]{babel} %Deutsche Silbentrennung
\usepackage[utf8]{inputenc} %Deutsche Umlaute
\usepackage{float}
\usepackage{graphicx}
\usepackage{minted}

\DeclareGraphicsExtensions{.pdf,.png,.jpg}

\makeatletter
\author{Vorkursteam der Fachschaft Informatik}
\let\Author\@author

% dark Mode
\ExplSyntaxOn
\RequirePackage{pagecolor,xcolor, graphicx} % Used for dark Mode
\bool_gset_false:N \g_dark_mode_bool % Disable by default
\newcommand{\enableDarkMode}{ %Command to enable Dark Mode (only works before \begin{document})
	\definecolor{anthrazitgrau}{HTML}{293133}
	\pagecolor{anthrazitgrau}
	\color{white}

	\cs_if_exist:NT \setbeamercolor {
		\setbeamercolor*{smallrule}{bg=.}
		\setbeamercolor*{normal~text}{bg=,fg=.}
		\setbeamercolor*{background canvas}{parent=normal~text}
		\setbeamercolor*{section~in~toc}{parent=normal~text}
		\setbeamercolor*{subsection~in~toc}{parent=normal~text,fg=\thepagecolor}
		\setbeamercolor*{footline}{parent=normal~text}
		\setbeamercolor{block~title~alerted}{fg=white,bg=white!20!\thepagecolor}
		\setbeamercolor*{block~body}{bg=white!10!\thepagecolor}
		\setbeamercolor*{block~body~alerted}{bg=\thepagecolor}
	}
	\cs_if_exist:NT \setbeamertemplate {
		\setbeamertemplate{subsection~in~toc~shaded}[default][50]
	}
	\bool_gset_true:N \g_dark_mode_bool
	% Prefer inverted Logo with dark Mode
	% \IfFileExists{tuda_logo_inverted.pdf}{\tl_gset:Nn \g_ptxcd_logofile_tl {tuda_logo_inverted.pdf}}{}
	% \hbox_gset:Nn \g__ptxcd_logo_box {% Update Logo Box
	% 	\makebox[2.2\c_ptxcd_logoheight_dim][l]{\includegraphics[height=\c_ptxcd_logoheight_dim]{\g_ptxcd_logofile_tl}}%
	% }
}

% dark Mode Makros
\prg_new_conditional:Nnn \__ptxcd_if_dark_mode: {T,F,TF} { % Conditional to check if dark Mode is active
	\bool_if:NTF \g_dark_mode_bool
	{\prg_return_true:}
	{\prg_return_false:}
}

\cs_set_eq:NN\IfDarkModeT \__ptxcd_if_dark_mode:T % Easy dark Mode check for use in document
\cs_set_eq:NN\IfDarkModeF \__ptxcd_if_dark_mode:F
\cs_set_eq:NN\IfDarkModeTF \__ptxcd_if_dark_mode:TF

\newcommand{\includeinvertablegraphics}[2][]{% Grafik wird beim Dark Mode automatisch Invertiert (rgb)
	\IfDarkModeTF{\includegraphics[decodearray={1.0~0.0~1.0~0.0~1.0~0.0},#1]{#2}}{\includegraphics[#1]{#2}}
}
\newcommand{\includeinvertablegrayscalegraphics}[2][]{% Grafik wird beim Dark Mode automatisch Invertiert (grayscale)
	\IfDarkModeTF{\includegraphics[decodearray={1.0 0.0},#1]{#2}}{\includegraphics[#1]{#2}}
}

% DARK_MODE environment check (enable if DARK_MODE=1)
\sys_get_shell:nnN { kpsewhich ~ --var-value ~ DARK_MODE } { } \l_dark_mode_env_var_tl
\tl_trim_spaces:N \l_dark_mode_env_var_tl
\tl_if_eq:NnT \l_dark_mode_env_var_tl {1} {\enableDarkMode{}}
\ExplSyntaxOff

% macros
\renewcommand{\arraystretch}{1.2} % Höhe einer Tabellenspalte minimal erhöhen
\newcommand{\N}{{\mathbb N}}
\newcommand{\code}{\inputminted[]{python}}
\newmintedfile[pythonfile]{python}{
	fontsize=\small,
	style=\IfDarkModeTF{fruity}{friendly},
	linenos=true,
	numberblanklines=true,
	tabsize=4,
	obeytabs=false,
	breaklines=true,
	autogobble=true,
	encoding="utf8",
	showspaces=false,
	xleftmargin=20pt,
	frame=single,
	framesep=5pt,
}
\newmintinline{python}{
	style=\IfDarkModeTF{fruity}{friendly},
	encoding="utf8"
}

\definecolor{codegray}{HTML}{eaf1ff}
\newminted[bashcode]{awk}{
	escapeinside=||,
	fontsize=\small,
	style=\IfDarkModeTF{fruity}{friendly},
	linenos=true,
	numberblanklines=true,
	tabsize=4,
	obeytabs=false,
	breaklines=true,
	autogobble=true,
	encoding="utf8",
	showspaces=false,
	xleftmargin=20pt,
	frame=single,
	framesep=5pt
}

\let\origpythonfile\pythonfile
\renewcommand{\pythonfile}[1]{\pythonfileh{#1}{}}
\newcommand{\pythonfileh}[2]{\origpythonfile[#2]{#1}}

\newcommand*{\ditto}{\texttt{\char`\"}}

%Includes
\usepackage{epstopdf}
\usepackage{wrapfig}
\usepackage{tipa}
\usepackage{tikz}
\usetikzlibrary{calc,shapes,arrows}
%tip: use http://l04.scarfboy.com/coding/phonetic-translation?from=ipa&fromtext=%CB%88pa%C9%AA%CE%B8n%CC%A9&to=tipa
%for converting ipa

\graphicspath{ {./media/} }

\def\shortyear#1{\expandafter\shortyearhelper#1}
\def\shortyearhelper#1#2#3#4{#3#4}

\newcount\NextYear
\NextYear = \year
\advance\NextYear by 1

\newcommand\NextYearShort{\shortyear{\the\NextYear}}

% notes
\usepackage{pgfpages}
\setbeamertemplate{note page}[plain]
%\setbeameroption{show notes on second screen}

% macro for change speaker sign
\newcommand{\changespeaker}{
	\begin{tikzpicture}[line width=.6mm, shorten >= 3pt, shorten <= 3pt]

		\coordinate (c1);
		\coordinate[right of=c1] (c2);

		\draw[rectangle, draw=red!80, fill=red!80, align=center, rounded corners] ($(c1.north west)+(0,-0.3)$) rectangle ($(c2.south east)+(0, 0.3)$) {};
		\draw[->,white] (c1)[bend left] to node[auto] {} (c2);
		\draw[->,white] (c2)[bend left] to node[auto] {} (c1);
	\end{tikzpicture}
}

%Listing-Style pyhon
\title[Programmiervorkurs]{Programmiervorkurs Wintersemester \the\year/\NextYearShort}
\subtitle{{\small der Fachschaft Informatik}}
\logo*{\includegraphics{../globalMedia/bildmarke_ohne_rand}}
\institute{Fachschaft Informatik}
\date{Wintersemester \the\year/\NextYearShort}


% macros
\newcommand{\livecoding}{
		\ifdefined\StreamSlides
			\begin{frame}
				\frametitle{\insertsectionhead \\  {\small \insertsubsectionhead}}\centering \huge 	\vskip 2cm\textbf{\textcolor{red}{Live-Coding}}
			\end{frame}
		\fi
	}

%\newcommand{\slidehead}{\frametitle{\insertsectionhead \\ {\small \insertsubsectionhead}}\vspace{3mm}}
\newcommand{\slidehead}{\frametitle{\insertsectionhead} \framesubtitle{\insertsubsectionhead}\vspace{3mm}}
\newcommand{\tocslide}{\begin{frame}[t]\frametitle{Inhaltsverzeichnis}\vspace{3mm}{\small\tableofcontents[subsectionstyle=shaded]}\end{frame}}

\newcommand{\nextvid}[2]{
	\ifdefined\StreamSlides
	\else
		\section{Nächstes Video}
		\begin{frame}[t]
			\slidehead
			\begin{block}{Nächstes Video}
				\vspace{0.5cm}
				#1
				\vspace{0.5cm}
			\end{block}
			\ifx\hfuzz#2\hfuzz
				\vspace{2.5cm}
			\else
				{\begin{block}{Bonus Video}
					\vspace{0.5cm}
					#2
					\vspace{0.5cm}
				\end{block}}
			\fi
			Danke fürs Zuschauen!\\
			Links zu den Folien und Quellen sind in der Videobeschreibung.
		\end{frame}
	\fi
}


\usepackage{verbatim}
\usetikzlibrary{decorations.pathreplacing}
\usetikzlibrary{shapes.misc}


% colors
\definecolor{lightpetrol}{RGB}{0,223,194}

% dark Mode
\ExplSyntaxOn
\RequirePackage{pagecolor,xcolor, graphicx} % Used for dark Mode
\bool_gset_false:N \g_dark_mode_bool % Disable by default
\newcommand{\enableDarkMode}{ %Command to enable Dark Mode (only works before \begin{document})
	\pagecolor{black}
	\color{white}
	\setbeamercolor*{smallrule}{bg=white}
	\setbeamercolor*{normal~text}{bg=,fg=white}
	\setbeamercolor*{background canvas}{parent=normal~text}
	\setbeamercolor*{section~in~toc}{parent=normal~text}
	\setbeamercolor*{subsection~in~toc}{parent=normal~text,fg=black}
	\setbeamertemplate{subsection~in~toc~shaded}[default][50]
	\setbeamercolor*{footline}{parent=normal~text}
	\setbeamercolor{block~title~alerted}{fg=white,bg=white!20!black}
	\setbeamercolor*{block~body}{bg=white!10!black}
	\setbeamercolor*{block~body~alerted}{bg=black}
	\bool_gset_true:N \g_dark_mode_bool
	\IfFileExists{tuda_logo_inverted.pdf}{\tl_gset:Nn \g_ptxcd_logofile_tl {tuda_logo_inverted.pdf}}{} % Prefer inverted Logo with dark Mode
	\hbox_gset:Nn \g__ptxcd_logo_box {% Update Logo Box
		\makebox[2.2\c_ptxcd_logoheight_dim][l]{\includegraphics[height=\c_ptxcd_logoheight_dim]{\g_ptxcd_logofile_tl}}%
	}
}

% dark Mode Makros
\prg_new_conditional:Nnn \__ptxcd_if_dark_mode: {T,F,TF} { % Conditional to check if dark Mode is active
	\bool_if:NTF \g_dark_mode_bool
	{\prg_return_true:}
	{\prg_return_false:}
}

\cs_set_eq:NN\IfDarkModeT \__ptxcd_if_dark_mode:T % Easy dark Mode check for use in document
\cs_set_eq:NN\IfDarkModeF \__ptxcd_if_dark_mode:F
\cs_set_eq:NN\IfDarkModeTF \__ptxcd_if_dark_mode:TF

\newcommand{\includeinvertablegraphics}[2][]{% Grafik wird beim Dark Mode automatisch Invertiert (rgb)
	\IfDarkModeTF{\includegraphics[decodearray={1.0~0.0~1.0~0.0~1.0~0.0},#1]{#2}}{\includegraphics[#1]{#2}}
}
\newcommand{\includeinvertablegrayscalegraphics}[2][]{% Grafik wird beim Dark Mode automatisch Invertiert (grayscale)
	\IfDarkModeTF{\includegraphics[decodearray={1.0 0.0},#1]{#2}}{\includegraphics[#1]{#2}}
}

% DARK_MODE environment check (enable if DARK_MODE=1)
\sys_get_shell:nnN { kpsewhich ~ --var-value ~ DARK_MODE } { } \l_dark_mode_env_var_tl
\tl_trim_spaces:N \l_dark_mode_env_var_tl
\tl_if_eq:NnT \l_dark_mode_env_var_tl {1} {\enableDarkMode{}}
\ExplSyntaxOff

% notes
\usepackage{pgfpages}
\setbeamertemplate{note page}[plain]
%\setbeameroption{show notes on second screen}

% macro for change speaker sign
\newcommand{\changespeaker}{
	\begin{tikzpicture}[line width=.6mm, shorten >= 3pt, shorten <= 3pt]

		\coordinate (c1);
		\coordinate[right of=c1] (c2);

		\draw[rectangle, draw=red!80, fill=red!80, align=center, rounded corners] ($(c1.north west)+(0,-0.3)$) rectangle ($(c2.south east)+(0, 0.3)$) {};
		\draw[->,white] (c1)[bend left] to node[auto] {} (c2);
		\draw[->,white] (c2)[bend left] to node[auto] {} (c1);
	\end{tikzpicture}
}

%Listing-Style pyhon
\title[Programmiervorkurs]{Programmiervorkurs Wintersemester \the\year/\NextYearShort}
\subtitle{{\small der Fachschaft Informatik}}
\logo*{\includegraphics{../globalMedia/bildmarke_ohne_rand}}
\institute{Fachschaft Informatik}
\date{Wintersemester \the\year/\NextYearShort}


% macros
\newcommand{\livecoding}{\begin{frame}\frametitle{\insertsectionhead \\  {\small \insertsubsectionhead}}\centering \huge \vskip 2cm\textbf{\textcolor{red}{Live-Coding}}\end{frame}}

%\newcommand{\slidehead}{\frametitle{\insertsectionhead \\ {\small \insertsubsectionhead}}\vspace{3mm}}
\newcommand{\slidehead}{\frametitle{\insertsectionhead} \framesubtitle{\insertsubsectionhead}\vspace{3mm}}
\newcommand{\tocslide}{\begin{frame}[t]\frametitle{Inhaltsverzeichnis}\vspace{3mm}{\small\tableofcontents[subsectionstyle=shaded]}\end{frame}}


% colors
\definecolor{lightpetrol}{RGB}{\IfDarkModeTF{0,136,119}{0,223,194}}

\begin{document}

%Deckblatt
\subtitle{Datentypen und Operatoren}
\maketitle

%Inhaltsverzeichnis
\tocslide

\section{Was sind Typen?}
\subsection{Überblick}
\begin{frame}
	\slidehead
	\begin{itemize}
		\item Daten können unterschiedliche \textit{Typen} haben
		\item Wichtige primitive Datentypen:
	\end{itemize}
	\begin{center}
		\begin{tabular}[h]{|c|c|c|}
			\hline
			\textbf{Typ}&\textbf{Wertebereich}&\textbf{Beispiele}\\
			\hline
			bool&Wahrheitswert& \pythoninline{True}, \pythoninline{False}\\
			\hline
			int&Ganzzahl& 42, -2 \\
			\hline
			float&Fließkommazahl& 42.5, 0.1\\
			\hline
			\end{tabular}\\
		\end{center}
		{\footnotesize Die Liste ist unvollständig, es gibt noch weitere Typen.} % Zum Beispiel gibt es einen für Komplexe Zahlen.
		\pause
		\begin{block}{Wichtig}
			Nutze einen Punkt als Dezimaltrenner für Fließkommazahlen!
		\end{block}
	\end{frame}

	\subsection{Technische Details}
	\begin{frame} % Technische Details die man sich nicht merken muss, die aber trotzdem Pflicht sind zu sagen.
		\slidehead
		Es ist nicht ganz so einfach...

		\begin{itemize}
			\item Bei anderen Sprachen muss der verwendete Typ explizit festgelegt werden % Hier kann ich kein Beispiel geben, weil wir noch keine Variablen eingeführt haben
		\end{itemize}

		In Python gilt:
		\begin{itemize}
			\item \pythoninline{float} kann nur Zahlen zwischen \pythoninline{-1.7976931348623157e+308} und \pythoninline{1.7976931348623157e+308} speichern
			\item Kleinster positiver \pythoninline{float} ist \pythoninline{2.2250738585072014e-308}
			\item \pythoninline{int} kann in Python beliebig große Zahlen speichern
			\begin{itemize}
				\item Das ist in den meisten Sprachen nicht so!
		\end{itemize}

		\end{itemize}
	\end{frame}

	\section{Operatoren}
	\begin{frame}
		\slidehead

		\begin{itemize}
			\item Symbol, das den \textbf{Interpreter} zur Ausführung einer Aktion veranlasst
			\item Operatortypen:
		\end{itemize}
		\tikzstyle{rect} = [rectangle, rounded corners, minimum width=3cm, minimum height=1cm,text centered, draw=black, fill=red!10]
		\centering
		\begin{tikzpicture}[node distance=1.5cm]
			\node (1) [rect, text width=5cm, minimum height=1cm] {\textbf{Mathematische Operatoren}};
			\node (2) [rect, below of=1, text width=5cm, minimum height=1cm] {\textbf{Vergleichsoperatoren}};
			\node (3) [rect, below of=2, text width=5cm, minimum height=1cm] {\textbf{Bitoperatoren}};
		\end{tikzpicture}
	\end{frame}

	\subsection{Mathematische Operatoren}
	\begin{frame}
		\slidehead
		\tikzstyle{rect} = [rectangle, rounded corners, minimum width=3cm, minimum height=1cm,text centered, draw=black, fill=red!10]
		\vspace{0.1cm}
		\centering
		\begin{tikzpicture}[node distance=1.2cm]
			\node (1) [rect, text width=3cm, minimum height=1cm] {Addition: $+$};
			\node (2) [rect, below of=1, text width=3cm, minimum height=1cm] {Subtraktion: $-$};
			\node (3) [rect, below of=2, text width=3cm, minimum height=1cm] {Multiplikation: $*$};
			\node (4) [rect, node distance=3.5cm, right of=1, text width=3cm, minimum height=1cm] {Division: $/$};
			\node (5) [rect, below of=4, text width=3cm, minimum height=1cm] {Exponentiation: $**$};
			\node (6) [rect, below of=5, text width=3cm, minimum height=1cm] {Modulo: $\%$};
		\end{tikzpicture}
		\pause
		\begin{block}{Hinweis}
			\begin{itemize}
				\item Modulo liefert den \textit{Rest einer Division}: $23 \% 4 = 3$
				\item Arithmetische Operationen zwischen Ganzzahlen und Fließkommazahlen resultieren in \textit{Fließkommazahlen}! \pythoninline{1 + 2.0 = 3.0} % Für die*den Programmieranfänger*innen mag es zwar total normal scheinen, dass man Kommazahlen mit Integers addieren kann, aber für den Computer ist das ein wichtiger Unterschied! Jetzt kommt der Moment wo ich alle überzeugen kann, dass der Computer wirklich komplett anders denkt als wir Menschen. Er hat ein komplett anderes Konzept sich 0.1 vorzustellen.

				%Umfrage, glaubt ihr Computer können rechnen?
				\item \pythoninline{0.1 + 0.2 =} \pause \pythoninline{0.30000000000000004} $\Rightarrow$ Fließkommazahlen sind nicht exakt
			\end{itemize}
		\end{block}
	\end{frame}


	\subsection{Vergleichsoperatoren}
	\begin{frame}
		\slidehead

		\begin{itemize}
			\item Vergleichsoperatoren werten zu einem Wert vom Typ \pythoninline{bool} aus
			\item Beispiel:
			\pythonfile{listings/vergleich_Listing.py}
		\end{itemize}
	\end{frame}


	\subsection{Vergleichsoperatoren - Übersicht}
	\begin{frame}
		\slidehead

		\begin{center}
			\vskip -5 pt
			\begingroup
			\fontsize{10pt}{11pt}\selectfont
			\begin{tabular}{|l|c|c|c|}
				\hline
				\textbf{Operator} & \textbf{Symbol} & \textbf{Beispiel} & \textbf{Rückgabewert} \\ \hline
				gleich 					& == 	& 100 == 50; 	& False \\
				&  		& 50 == 50; 	& True \\ \hline
				ungleich 				& != 	& 100 != 50; 	& True \\
				&  		& 50 != 50 		& False \\ \hline
				größer als 				& > 	& 100 > 50; 	& True \\
				&  		& 50 > 50; 		& False \\ \hline
				größer als oder gleich 	& >= 	& 100 >= 50; 	& True \\
				& 		& 50 >= 50 		& True \\ \hline
				kleiner als 			& < 	& 100 < 50; 	& False \\
				& 		& 50 < 50; 		& False \\ \hline
				kleiner als oder gleich & <= 	& 100 <= 50; 	& False \\
				&  		& 50 <= 50; 	& True \\ \hline
			\end{tabular}
			\endgroup
		\end{center}

		\note{
		== geht auch für Strings!
		}
	\end{frame}

	\subsection{Bitoperatoren}
	\begin{frame}
		\slidehead
		Ja, es gibt auch Bitoperatoren ...
	\end{frame}

	\subsection{}
	\livecoding

	\subsection{Operatorpräzedenz}
	\begin{frame}
		\slidehead

		\pythonfile{listings/Rangfolge_Listing.py}
		\pause
		\pythonfile{listings/Rangfolge_Listing_parenthesis.py}

		\begin{itemize}
			\item Die Operatorpräzedenz gibt die \textbf{Reihenfolge} vor, in der Ausdrücke ausgewertet werden:
			\begin{enumerate}
				\item Arithmetische Ausdrücke
				\item Vergleiche
				\item Logische Ausdrücke (hierzu später mehr)
				\item Zuweisungen
			\end{enumerate}
		\end{itemize}
	\end{frame}

	\section{Zeichenketten}

	\begin{frame}
		\slidehead
		\begin{itemize}
			\item Habt ihr vorhin kennengelernt
			\pause
			\item Zeichenketten werden als \textit{Strings} bezeichnet
			\item Ein String kann beliebig viele Zeichen enthalten
			\pause
			\item Strings können mit \textbf{+} konkateniert werden
		\end{itemize}
	\end{frame}


	%Vielleicht später mal nützlich:

	%\begin{frame}
	%	\slidehead
	%	\begin{itemize}
	%		\item String Slicing (Zerschneiden von Text)
	%			\pythonfile{listings/String_Slicing_Listing.py}
	%		\begin{block}{Wichtig}
	%			\begin{itemize}
	%				\item \textbf{[von : bis]} ist \textbf{inklusive} "`von"' und \textbf{exklusive} "`bis"'
	%				\item Mathematisch: Ein Interval $[von, bis)$ bzw. $[von, bis[$
	%				\item Der erste Index beginnt mit der Zahl \textbf{0}
	%			\end{itemize}
	%		\end{block}
	%	\end{itemize}
	%\end{frame}

	%\begin{frame}
	%	\slidehead
	%	\begin{itemize}
	%		\item String Slicing Teil 2
	%			\pythonfile{listings/String_Slicing_Listing2.py}
	%		\begin{block}{Wichtig}
	%			\begin{itemize}
	%				\item \textbf{[: bis]} ist \textbf{exklusive} und \textbf{[von :]} ist \textbf{inklusive} der Grenze
	%				\item Auch hier beginnt der erste Index mit der Nummer \textbf{0}
	%				\item Es können auch beide Grenzen weggelassen werden
	%			\end{itemize}
	%		\end{block}
	%	\end{itemize}
	%\end{frame}

	\subsection{Funktionen und Operatoren}
	\begin{frame}
		\slidehead

		\begin{itemize}
			\item Länge eines Strings
			\pythonfile{listings/String_len_Listing.py}

			\item Umwandeln in Kleinbuchstaben
			\pythonfile{listings/String_lower_Listing.py}

			\item Umwandeln in Großbuchstaben
			\pythonfile{listings/String_upper_Listing.py}
		\end{itemize}
	\end{frame}


	\subsection{Konvertierung}
	\begin{frame}
		\slidehead
		\begin{block}{Warnung}
			Strings sind keine Zahlen
		\end{block}

		\pythonfile{listings/String_is_not_a_number_Listing.py}

		\pause

		\begin{itemize}
			\item Strings müssen zum Rechnen konvertiert werden
			\item \textbf{int(x)} konvertiert x zu int
			\item \textbf{float(x)} konvertiert x zu float
%			\item \textbf{bool(x)} konvertiert x zu boolean  % Das lassen wir weg, weil es durch leere Zeichenketten und Zahlen zu mehr Verwirrung sorgt als es hilfreich ist
			\item \textbf{str(x)} konvertiert x zu String
		\end{itemize}

	\end{frame}

	\begin{frame}
		\slidehead
		\begin{block}{Beispiel}
			\pythonfile{listings/String_conversion_Listing.py}
		\end{block}
	\end{frame}

	\livecoding

	\section{Quiz}
	\begin{frame}
		\slidehead

		\begin{itemize}
			\item Wie kann ich \textbf{'} in Strings verwenden?
			\pause
			\item Wo liegt der Unterschied zwischen \pythoninline{int} und \pythoninline{float}?
			\pause
			\item In welcher Reihenfolge wird \pythoninline{2+4*2 == 9} ausgewertet?
			\pause
			\item Was ist das Ergebnis?
			\pause
			\item Wofür gibt es Kommentare?
		\end{itemize}
	\end{frame}

	\end{document}
